\section{A normal Euler field: a complete constructive criterion}
Let $R=k[Y_1,\ldots,Y_n]$, let $\gamma:\A^1\hookrightarrow\A^n$ be a closed embedding,
and put $I=I_\gamma$, $B=R[t,I/t]$. For a derivation $\delta$, consider
the conditions
\begin{equation}\label{eq:normal}
\delta(R)\subset I,\qquad (\delta-\mathrm{id})(I)\subset I^2.
\end{equation}
They say that the field vanishes on the curve and acts as the identity on
first-order normal functions. Then
\[
E=\partial_t+t^{-1}\delta
\]
preserves $B$, since $E(Y_i)=\delta Y_i/t\in B$ and
$E(F/t)=(\delta F-F)/t^2\in B$ for $F\in I$.
Local finiteness of $\delta$ means that
$\Span_k\{\delta^jF:j\ge0\}$ is finite-dimensional for every $F\in R$.

\begin{theorem}\label{thm:euler}
For a fixed $\delta$ satisfying \eqref{eq:normal}, the following are equivalent:
\begin{enumerate}
\item $E$ is an LND on $B$;
\item $\delta$ is locally finite on $R$;
\item there are coordinates $R=k[p,q_1,\ldots,q_{n-1}]$ such that
$p\gamma=u$, $I=(q_1,\ldots,q_{n-1})$, and
$\delta=\sum_jq_j\partial_{q_j}$.
\end{enumerate}
Consequently, the existence of such a locally finite field is equivalent
to rectifiability of $\gamma$ over $k$.
\end{theorem}

\begin{proof}
In the localization, direct induction gives
\begin{equation}\label{eq:falling}
E^mF=t^{-m}\prod_{j=0}^{m-1}(\delta-j)F,
\qquad
E^m(F/t)=t^{-m-1}\prod_{j=1}^{m}(\delta-j)F.
\end{equation}
If $E$ is locally nilpotent, every $F$ is annihilated by a nonzero
polynomial in $\delta$, which gives local finiteness.

The Leibniz rule and \eqref{eq:normal} give
\[
(\delta-m)I^m\subset I^{m+1}\quad(m\ge0).
\]
Now suppose that $\delta$ is locally finite.
Let $V$ be a finite-dimensional $\delta$-stable space,
and let $M$ be the largest ordinary degree of its elements.
Since $I\subset\mathfrak m_{\gamma(0)}$, we have $V\cap I^{M+1}=0$.
Thus $\delta(\delta-1)\cdots(\delta-M)$ annihilates $V$.
Its roots are distinct and belong to $k$: semisimplicity and nonnegative
integer weights follow automatically, without passage to the
algebraic closure. We obtain
\[
R=\bigoplus_{d\ge0}R_d,\qquad \delta|_{R_d}=d.
\]
For a nonzero eigenvector, its exact $I$-adic order equals its weight,
because the induced action on $I^m/I^{m+1}$ is multiplication by $m$.
The ideals $I^m$ are graded: the spectral projectors are
polynomials in $\delta$. Hence
\[
I^m=\bigoplus_{d\ge m}R_d,\quad R_0\simeq R/I\simeq k[u],\quad
R_1\simeq I/I^2,\quad R=R_0[R_1].
\]
The last equality follows by induction on the weight: for $d\ge2$ we have
$R_d\subset I^2$, and each homogeneous factor in a product
has a smaller positive weight.

Freeness of $I/I^2$ does not require homological algebra.
Choose $P\gamma=u$, and put $v=\gamma'$, $a=\nabla P(\gamma)$.
Then $av=1$, and the row module
$L=\{\ell\in k[u]^n:\ell v=0\}$ is free of rank $n-1$
over the PID $k[u]$. The gradient defines an isomorphism
\[
I/I^2\longrightarrow L,\qquad F\longmapsto\nabla F(\gamma).
\]
An explicit lift of a row is
\[
F_\ell(Y)=\sum_i\ell_i(P(Y))\bigl(Y_i-f_i(P(Y))\bigr).
\]
Injectivity follows from Taylor expansion: a function whose value
and first gradient both vanish on the curve belongs to $I^2$.
Choose an invariant lift $p$ of the parameter $u$ and a basis $q_j$ of
$R_1$ over $R_0=k[p]$. Then $R=k[p,q_1,\ldots,q_{n-1}]$;
algebraic independence follows from transcendence degree $n$.
This gives the normal form of $\delta$.
Finally, $B=k[t,p,q_1/t,\ldots,q_{n-1}/t]$, and $E=\partial_t$
in these coordinates. This proves all the implications.
\end{proof}

\subsection{Finite verification and reconstruction of the automorphism}
It suffices to exhibit the exact stability condition
\[
V_N=\Span_k\{1,\delta^jY_i:0\le j\le N, 1\le i\le n\},\qquad
\delta^{N+1}Y_i\in V_N.
\]
Spaces spanned by products of elements of $V_N$ are stable by the Leibniz rule,
so $\delta$ is locally finite on all of $R$.
If the degrees of elements of $V_N$ are at most $M$, the spectral
projection of a polynomial $F$ onto weight $j$ is computed by
\[
\pr_j(F)=\prod_{\substack{0\le i\le L\\i\ne j}}
\frac{\delta-i}{j-i}F,\qquad L\ge M\deg F.
\]
Thus
$p=\pr_0(P)$, $q_j=\pr_1(F_{\ell_j})$ are the required coordinates.
The inverse map is reconstructed from the weight components of
$Y_i$ and finitely many transverse derivatives. The detailed formula is retained
in \file{euler/theory.md}. No universal a priori bound on $N$ is given here.

A positive control reconstructs the nonlinear coordinates
\[
r=z-x^2,\qquad q=y-r^2,\qquad p=x-q^2.
\]
The corresponding space $V_4$ has dimension $9$ and is exactly
stable. This checks the implementation; it does not replace the general proof.

\subsection{Necessary conditions that are still insufficient}
The normal form implies
\[
\Div\delta=n-1,\qquad (\delta Y_1,\ldots,\delta Y_n)=I.
\]
The equality of ideals describes the entire fixed-point scheme, rather than
only its set of points. Constancy of the divergence of an LF field is a classical
fact \cite{Nowicki}; no priority for this fact is claimed.

Even both conditions together with \eqref{eq:normal} are insufficient.
For the line $I=(y,z)$, the field
\[
\delta=yz\partial_x+y\partial_y+(z+y^2x)\partial_z
\]
passes all these tests. However, $\delta^{2m}x$ contains
$xy^{3m}$ with a positive coefficient. It arises by repeating
the path $x\mapsto yz\mapsto y^3x$; all coefficients are nonnegative,
so no cancellation occurs. The degrees are unbounded, hence the field is not LF.
Local finiteness of the normal field must not be confused with local
nilpotence: in the correct normal form $\delta^m q_j=q_j$, while $E$
is the locally nilpotent derivation.

\section{The exact limitation of the naive normal field}
For any parameter polynomial $P\gamma=u$, put
\[
\delta_P=\sum_i\bigl(Y_i-f_i(P)\bigr)\partial_{Y_i}.
\]
This field always satisfies \eqref{eq:normal}, but has the following
essential limitation.

\begin{theorem}
There exists $P\gamma=u$ for which $\delta_P$ is locally finite
if and only if there exists an affine-linear $L$ such that
$L\gamma=u$.
\end{theorem}
\begin{proof}
If $P=L$ is affine-linear, then $\delta_P(P)=0$ and
$\delta_P^2Y_i=\delta_PY_i$, giving local finiteness on the generators.
Conversely, the necessary divergence identity is
\[
\sum_i f_i'(P)\partial_iP=1.
\]
By an affine change of coordinates, make the positive degrees of the
nonconstant $f_i$ pairwise distinct; the remaining components can be made zero.
The form of $\delta_P$ is preserved under this change. If there is no linear
parameter, all positive degrees are at least $2$.
Let $d=\deg P$, and let $j$ correspond to the largest degree of an $f_j$
with $\partial_jP\ne0$. The degree of this summand is at least
$(\deg f_j-1)d$, whereas each preceding summand has degree at most
$(\deg f_i)d-1\le(\deg f_j-1)d-1$.
The leading term cannot cancel, so the sum cannot equal $1$.
If there are no active nonzero components, the sum is zero,
which is also impossible.
\end{proof}

For harmonic models with distinct indices $a,b,c\ge2$, there is no
linear parameter. Thus \emph{no} choice of $P$, of any degree,
makes this field LF. The exclusion also applies to the already rectified
$H(3,4,5)$ and $H(3,5,7)$, so it is not an obstruction to rectifiability
of the curve itself.

\section{General normal fields: exact spaces of candidates}
We seek all $\delta=a_1\partial_x+a_2\partial_y+a_3\partial_z$ with
$\deg a_i\le d$, satisfying \eqref{eq:normal} and $\Div\delta=2$.
Each $a_i$ is expanded in a full basis of $I\cap R_{\le d}$.
Writing $v=\gamma'$, the second normal condition is equivalent to
the columns of $J_a(\gamma)-\mathrm{Id}$ being multiples of the unimodular $v$.
This gives a rational linear system $Ac=b$ covering
all fields of the specified degree bound.

\begin{center}\small
\begin{tabular}{@{}lrrrrl@{}}\toprule
Model & $d$ & $\dim I_{\le d}$ & Unknowns & $\operatorname{rank}A$ & Result\\\midrule
$H(3,7,8)$ &7&63&189&189&inconsistent\\
$H(3,7,8)$ &8&100&300&276&family $24$\\
$H(4,5,7)$ &6&42&126&---&inconsistent\\
$H(4,5,7)$ &7&71&213&194&family $19$\\\bottomrule
\end{tabular}
\end{center}
Inconsistency is certified by an exact rational vector satisfying
$\lambda^{\mathsf T}A=0$, $\lambda^{\mathsf T}b=1$; the conclusion therefore
persists under any characteristic-zero field extension.
For the consistent systems, a rational point and a full basis of
directions are provided. Exact null vectors give upper bounds on rank;
matching modular lower bounds certify completeness.

Thus the first degrees at which the \emph{normal conditions together
with global divergence $2$} are consistent are $8$ and $7$.
These are not minimum degrees of LF fields or rectifiers.
For $H(4,5,7)$, degree $6$ is already impossible without the divergence condition.
For $H(3,7,8)$, the normal conditions alone at degree $7$ have
an $11$-dimensional family; it is global divergence that excludes it.

\subsection{Two selected representatives are rigorously excluded}
In the $H(3,7,8)$ family, the selected rational representative preserves
$(x,y)$ and acts on the quotient $k[z]$ as $\psi(z)\partial_z$,
with $\deg\psi=6$. In the $H(4,5,7)$ family, the selected representative preserves
$(y,z)$ and induces $\varphi(x)\partial_x$, with $\deg\varphi=7$.
For a polynomial $a$ of degree $e\ge2$ in characteristic zero,
\[
\deg((a\partial_v)^m v)=1+(e-1)m,\qquad
\operatorname{lc}((a\partial_v)^m v)
=\operatorname{lc}(a)^m\prod_{j=0}^{m-1}(1+(e-1)j)\ne0.
\]
Hence neither representative is LF, and the corresponding $E$ are not LNDs.
The full coefficients and separate checks are included in the JSON certificates.
The axes are not preserved by the entire families: respectively $13$ of $24$ and
$11$ of $19$ selected basis directions fail to preserve them.
These counts concern the chosen bases; they are not dimensions of
excluded subfamilies. No exclusion of all complex parameter values is established.

\subsection{Real topology and conjugation}
By \cite[p.~2, Fig.~3]{KP}, $H(3,7,8)$ has type $T(2,5)$
up to mirror image, whereas $H(4,5,7)$ has type $5_2$.
Consequently, these models admit no real polynomial rectifiers.
By Theorem~\ref{thm:euler}, \emph{every} suitable LF field for them
must have nonreal coefficients. This is a known topological
restriction on the field of definition, not a complex obstruction.

If a suitable complex $\delta$ commuted with $\bar\delta$,
their average would be a real LF field satisfying the same normal
conditions. Commuting LF operators have finite-dimensional joint
cyclic spaces, so the average is indeed LF.
This would give the forbidden real rectification. Similarly,
commuting $E$ and $\bar E$ would give a real LND with slice $t$.
Therefore, both commutators are nonzero in these cases. Moreover,
\[
[E,\bar E]=t^{-2}\bigl([\delta,\bar\delta]+\delta-\bar\delta\bigr),
\]
so the two commutation conditions must not be identified.
